\section{Phát biểu bài toán, mục tiêu và phạm vi}
\label{sec:bai-toan}
\subsection{Đối tượng sử dụng và vai trò hỗ trợ}
Chatbot được thiết kế để hỗ trợ cán bộ, giảng viên tại Đại học Bách khoa
Hà Nội tìm hiểu và chuẩn bị các thủ tục thuộc phạm vi đề tài; đồng thời
hỗ trợ chuyên viên tra cứu căn cứ khi giải đáp những trường hợp này.
Các vai trò được phân thành hai nhóm: sử dụng trực tiếp và chuẩn bị nội
dung, vận hành hệ thống.

\Needspace{6\baselineskip}
\NhomTuVan{Đối tượng sử dụng trực tiếp}
\begin{itemize}[leftmargin=*,itemsep=2pt,topsep=3pt,before={\interlinepenalty=10000}]
\item \textbf{Cán bộ, giảng viên HUST:} mô tả công việc cần thực hiện và cung
 cấp dữ kiện liên quan để nhận hướng dẫn về điều kiện, tài liệu cần chuẩn
 bị và đầu mối liên hệ.
\item \textbf{Chuyên viên tại Ban Khoa học--Công nghệ, Ban Tài chính--Kế
 hoạch, Ban Tổ chức--Nhân sự và các đơn vị liên quan:} tra cứu quy định,
 kiểm tra căn cứ và tham khảo hướng dẫn khi giải đáp một trường hợp thuộc
 phạm vi chatbot.
\end{itemize}

\Needspace{6\baselineskip}
\NhomTuVan{Vai trò chuẩn bị nội dung và vận hành}
\begin{itemize}[leftmargin=*,itemsep=2pt,topsep=3pt,before={\interlinepenalty=10000}]
\item \textbf{Chuyên viên phụ trách nội dung nghiệp vụ:} kiểm tra văn bản
 nguồn, phạm vi áp dụng, điều kiện và đáp án tư vấn thuộc lĩnh vực phụ
 trách; xác nhận nội dung trước khi đưa vào sử dụng.
\item \textbf{Quản trị viên hệ thống:} cập nhật tài liệu và nội dung đã được
 xác nhận, đồng thời quản lý quyền truy cập theo phạm vi sử dụng được phép.
\end{itemize}
Trong thiết kế hệ thống, một cá nhân có thể đảm nhiệm nhiều vai trò.
Chuyên viên phụ trách nội dung và quản trị viên được sử dụng chức năng
hội thoại với tư cách người hỏi. Hoạt động nhận tư vấn được phân biệt
với trách nhiệm xác nhận nghiệp vụ và vận hành hệ thống.

Đề tài không phục vụ cá nhân ngoài HUST hoặc sinh viên không đồng thời
đảm nhiệm vai trò cán bộ, giảng viên hay chuyên viên. Cán bộ đang học sau
đại học được tư vấn theo vai trò công tác và thủ tục liên quan.

Quyền sử dụng chức năng và truy cập tài liệu được thiết kế theo vai trò
và phạm vi nguồn được phép sử dụng. Việc thuộc nhóm đối tượng phục vụ
không mặc nhiên cấp quyền khai thác toàn bộ tài liệu nội bộ.

\subsection{Phát biểu bài toán}
Bài toán đặt ra là từ yêu cầu về một thủ tục và thông tin của người hỏi,
hệ thống hình thành hướng dẫn phù hợp với trường hợp đó. Trong báo cáo,
\emph{tình huống} được hiểu là công việc cần thực hiện cùng những dữ kiện
có ảnh hưởng đến hướng dẫn. Bài toán gồm các thành phần sau:

\textbf{Đầu vào} là câu hỏi tiếng Việt, dữ kiện người dùng cung cấp, nội
dung đã trao đổi và tập nguồn được phép sử dụng. Dữ kiện có thể gồm loại
nhiệm vụ, vai trò của người hỏi, đơn vị, nguồn kinh phí và thời điểm. Chỉ
những thông tin ảnh hưởng đến câu trả lời mới cần được làm rõ; thông tin
người dùng khai báo chưa được coi là dữ kiện đã xác minh.

\textbf{Nhiệm vụ cần giải quyết} là xác định quy định liên quan, làm rõ
dữ kiện quyết định và đối chiếu điều kiện áp dụng cho trường hợp đang hỏi.
Nếu một yêu cầu liên quan đến nhiều nội dung, từng phần phải được xét trên
cùng thông tin của tình huống.

\textbf{Đầu ra} là hướng dẫn về điều kiện, tài liệu hoặc biểu mẫu cần
chuẩn bị, mốc cần lưu ý và bước hoặc đầu mối tiếp theo. Hướng dẫn kèm dữ
kiện đang sử dụng và đoạn nguồn hỗ trợ. Khi chưa đủ thông tin hoặc căn cứ,
hệ thống trả lời phần đã rõ, đồng thời nêu câu hỏi bổ sung hoặc phần chưa
thể kết luận.

Với ví dụ hội thảo ở Mục~\ref{sec:boi-canh}, đầu vào ban đầu là dự định
tham dự bằng kinh phí đề tài. Sau khi làm rõ các dữ kiện cần thiết, đầu ra
cần chỉ ra quy định áp dụng, tài liệu cần chuẩn bị và đầu mối liên hệ.
Nếu người hỏi sửa nguồn kinh phí hoặc thời điểm, hướng dẫn liên quan phải
được cập nhật và chỉ rõ phần thay đổi.

\subsection{Mục tiêu và sản phẩm của đề tài}
Mục tiêu sử dụng là giúp người hỏi biết mình cần đáp ứng điều kiện nào,
chuẩn bị những gì và thực hiện bước nào tiếp theo. Người hỏi cần kiểm tra
được căn cứ của hướng dẫn và nhận biết phần còn phải trao đổi với đơn vị
phụ trách.

Các sản phẩm cần xây dựng gồm ứng dụng chatbot tiếng Việt, kho tri thức
cho năm nhóm thủ tục và bộ tình huống kiểm thử có nguồn, đáp án đối chiếu.
Những yêu cầu về sự phù hợp của hướng dẫn với dữ kiện, khả năng truy vết
căn cứ và tính nhất quán qua hội thoại được cụ thể hóa thành F01--F05 tại
Mục~\ref{sec:yeu-cau}.

\subsection{Phạm vi tư vấn}
Ba miền nội dung của đề tài tương ứng với Ban Khoa học--Công nghệ (KHCN),
Ban Tài chính--Kế hoạch (TCKH) và Ban Tổ chức--Nhân sự (TCNS) tại HUST.
Phạm vi được giới hạn ở năm nhóm thủ tục trong Bảng~\ref{tab:pham-vi}, gồm
các nhánh tài chính và nhân sự liên quan trực tiếp. Năm nhóm này không bao
phủ toàn bộ công việc của ba Ban.

\begin{table}[H]
\centering\small
\begin{tabularx}{\textwidth}{|P{28mm}|Y|Y|}
\hline
\textbf{Nhóm thủ tục} & \textbf{Câu hỏi đại diện} & \textbf{Nội dung hướng dẫn} \\\hline
T1\newline Đăng ký nhiệm vụ & Tôi cần đáp ứng điều kiện gì và chuẩn bị những tài liệu nào để đăng ký nhiệm vụ khoa học và công nghệ? & Đối tượng, điều kiện theo loại nhiệm vụ; đợt đăng ký, tài liệu, biểu mẫu và đầu mối. \\ \hline
T3\newline Điều chỉnh, gia hạn & Khi cần thay đổi tiến độ, nội dung, chủ nhiệm hoặc kinh phí, tôi phải lưu ý điều kiện và mốc đề nghị nào? & Trường hợp cần điều chỉnh; điều kiện, thời hạn đề nghị và tài liệu theo nội dung thay đổi. \\ \hline
T4\newline Nghiệm thu & Trước khi đề nghị nghiệm thu, tôi cần chuẩn bị sản phẩm, báo cáo và tài liệu gì? & Yêu cầu sản phẩm, báo cáo, thời hạn và thành phần tài liệu phục vụ nghiệm thu. \\ \hline
T6\newline Hội thảo, công tác & Với hoạt động và nguồn kinh phí đang sử dụng, tôi cần đối chiếu quy định và chuẩn bị chứng từ nào? & Quy định nhiệm vụ, nguồn và phạm vi định mức; chứng từ, thời hạn tạm ứng hoặc quyết toán liên quan. \\ \hline
T10\newline Cử đi đào tạo, nghiên cứu & Tôi cần đáp ứng điều kiện, chuẩn bị tài liệu và thực hiện nghĩa vụ nào liên quan đến việc cử đi? & Điều kiện, thời gian, tài liệu, nghĩa vụ và hỗ trợ liên quan đến việc cử đi. \\
\hline
\end{tabularx}
\caption{Phạm vi năm nhóm nội dung tư vấn}
\label{tab:pham-vi}
\end{table}

Một câu hỏi có thể cần nguồn thuộc một hoặc nhiều Ban. Chẳng hạn, cán bộ
được cử đi đào tạo khi đang chủ trì nhiệm vụ có thể cần xem cả điều kiện
cử đi, trách nhiệm với nhiệm vụ và phần hỗ trợ tài chính được hỏi. Các
nguồn được đối chiếu theo nội dung liên quan; điều này không xác định một
tuyến luân chuyển hồ sơ qua các Ban.

\subsection{Giới hạn hỗ trợ và điều kiện áp dụng}
Chatbot hướng dẫn người dùng chuẩn bị thủ tục; không tiếp nhận, thẩm định,
luân chuyển, ký, nộp hoặc duyệt hồ sơ. Hệ thống không tính hay xác nhận
khoản chi thực tế, thực hiện thanh toán hoặc ghi dữ liệu vào hệ thống
nghiệp vụ của HUST. Việc liệt kê tài liệu và biểu mẫu trong câu trả lời
không xác nhận bộ hồ sơ của người hỏi đã hợp lệ.

Phạm vi không gồm tuyển dụng, tính kNC, đánh giá hoặc xếp loại cá nhân,
tra cứu hồ sơ nhân sự hay sinh viên. Sở hữu trí tuệ và hợp tác nghiên cứu
không thuộc phạm vi đề tài.
Giới hạn tra cứu hồ sơ cá nhân áp dụng cho mọi người sử dụng, kể cả cán
bộ, giảng viên và chuyên viên thuộc nhóm được phục vụ.

Nguồn tư vấn phải được kiểm tra và được phép sử dụng, đúng đối tượng, đơn
vị, loại nhiệm vụ và phiên bản áp dụng. Nguồn lịch sử được bổ sung khi nhiệm
vụ đang thực hiện, nghĩa vụ kéo dài hoặc giấy tờ tồn đọng cần xét chuyển
tiếp; không cam kết tra cứu mọi sự việc đã kết thúc từ một năm tùy chọn.
Nhánh còn thiếu nguồn phải được ghi nhận là chưa hoàn thành, kèm giới hạn
trả lời; không tự đặt quy tắc hoặc loại nhánh đó khỏi phạm vi đã chốt.

% Kết thúc phần xác định phạm vi trước khi chuyển sang phân tích yêu cầu.
\clearpage
